\section{Electronic response libraries}\label{app:E.16}

A spin-dependent coupling probes the spin density of the crystal, a velocity-dependent one its current, a scalar coupled to the electron's mass its kinetic-energy density, and each of these is screened differently from the charge density (\cref{sec:10.4,sec:10.7}). A rate for any of them needs a response function that most published calculations do not export: the full matrix of the independent-particle response, from which the screened responses to every operator follow by one Dyson equation. A library of these matrices, and the calculations that use it, follow the methods of \cref{app:E.6,app:E.7}.

\subsection{The response library}

For a crystal symmetric under the combination of inversion and time reversal, such as silicon or germanium, the response to the density $n_e(\vect{q})$ and the two transverse components of the current $j_{\mathrm{T}}(\vect{q})$ is fixed at each $(\vect{q},\omega)$ by 12 real functions (\cref{sec:10.4}). Six belong to the screened singlet block: the density structure factor $S_{nn}$, two density--current cross terms, and three symmetric current--current terms. Six more give the same block for independent particles, which also fixes the independent-particle triplet block, a quarter of the singlet one. The library holds these 12 functions on a grid of $(|\vect{q}|,\hat{\vect{q}},\omega)$ and, in addition, the full independent-particle matrix $R^0_{\vect{G}\vect{G}'}(\vect{q},\omega)$, so that other kernels can be applied later without repeating the band calculation. For operator $\mathcal{O}_2$, whose velocity squared brings in the kinetic-energy density $\mathcal{T}_e$, the library adds four functions in each version, the kinetic rows below, so that it holds 20 functions in all (\cref{sec:21.6}). On the truncation label $\mathfrak{T}$ (\cref{sec:8.5}) it sits at $n_B = 0$, $n_g = 0$, with $\mathcal{X}$ = IP and $\mathcal{X}$ = RPA with local fields.

For silicon the ground state is the shared setting of \cref{app:E.17} (PBE, SG15 at \SI{60}{Ry}, measured lattice constant) with a scissor to the measured gap, as in \cref{app:E.7}. A scissor keeps the band order of PBE, which in germanium is wrong (\cref{app:E.1}), so the germanium library is built only on germanium bands that pass the check of \cref{app:E.14} (valley order, both gaps, masses at L and $\mathrm{\Gamma}$), with quasiparticle $GW$ bands as the next route under that check; until such bands exist, no germanium version is released. Density vertices $\rho_{cv}(\vect{k},\vect{q}+\vect{G})$ come from matching Miller indices. The current vertex $\langle c|\tfrac12\{\vect{p},\ee^{\ii\vect{q}\cdot\vect{r}}\}|v\rangle/m_e$ needs the correction for the nonlocal pseudopotential of \cref{eq:E.7-vNL}; without it the transverse current of silicon is off by about \SI{10}{\percent}. The Dyson equation of \cref{eq:E.6-dyson} is solved on the joint index set of the density at $\vect{q}+\vect{G}$, with $|\vect{q}+\vect{G}| \le 8\,\alpha m_e$, and the current, which enters only as an external row and column. The grids follow published crystal form-factor calculations: $10^3$ $\vect{k}$ points without symmetry, $|\vect{q}| \le 25\,\alpha m_e$ in steps of $0.1\,\alpha m_e$, 30 directions in the irreducible wedge, $\omega \le \SI{50}{eV}$ in steps of \SI{0.05}{eV}, and 60 bands. The tables are converged when $S_{nn}$ at \SI{1}{keV} changes by less than \SI{3}{\percent} from $8^3$ to $12^3$ $\vect{k}$ points and by less than \SI{1}{\percent} when the local-field cutoff goes from 6 to $8\,\alpha m_e$. With the nonlocal term included, the $f$-sum rule of \cref{eq:E.7-fsum} has to hold to \SI{2}{\percent} below the band cutoff, the tolerance at which \cref{sec:21.6} releases a version and on which the transfer rule of \cref{we:22-fsum} relies; where the spectrum ends early, the fraction it holds is stored as a value per $|\vect{q}|$. Symmetry gives a check that costs nothing: in such a crystal the imaginary part of the density--current cross term and the antisymmetric part of the current--current block vanish, and numerically they stay below $10^{-3}$ of the real parts.

Storing the full matrix is impossible. At $|\vect{q}+\vect{G}| \le 8\,\alpha m_e$ (\SI{15}{\angstrom^{-1}}) it has about 2300 reciprocal-lattice vectors, \SI{85}{MB} per $(\vect{q},\omega)$, and on the full grid of $7.5\times10^6$ points it would fill some \SI{600}{TB}. The 12 functions need only the screened column at $\vect{G} = 0$ and the current rows, so the Dyson equation is solved for those few right-hand sides by an iterative solver, about $10^9$ operations per point, and never by inverting the matrix; the full $R^0_{\vect{G}\vect{G}'}$ is stored only on a coarse subset (25 values of $|\vect{q}|$, 3 directions, 100 frequencies) at a cutoff of $4\,\alpha m_e$, about \SI{10}{GB}. Building the bubble then dominates, about $10^{16}$ operations, some $10^3$ core-hours for silicon and four times that for germanium with spin--orbit coupling. The output is a set of versioned tables, which \cref{sec:21.6} uses as the reference electronic response.

The library also gives the screening of the current by local fields and the spin channel beyond independent particles. With the Dyson equation solved on the joint index set, the screened current--current response differs from its independent-particle value by a term built from the mixed blocks $R^0_{j_{\mathrm{T}},n(\vect{q}+\vect{G})}$, the density response to a transverse current. At $\vect{G} = 0$ this block vanishes when $\hat{\vect{q}}$ lies on a three- or fourfold axis, where a transverse current cannot drive a longitudinal density, so the screening of the current is purely a local-field effect, estimated at 10 to \SI{30}{\percent}. Its value per $(|\vect{q}|,\omega)$ is what settles the in-medium factor for the velocity-dependent operator in \cref{sec:10.4}. The spin channel has a Dyson equation that uses the matrix $R^0_{\vect{G}\vect{G}'}$ with a kernel made of the triplet exchange--correlation kernel and the screened electron--hole ladder, but without the exchange term that the charge channel carries. Since the exchange matrix element of a Wannier--Mott exciton is of relative order $(a/a_x)^3 \approx 10^{-3}$ in silicon, apart from its long-range part, which is the external screening and is absent from the spin channel by construction, the triplet ladder equals the singlet one to that accuracy. The ladder is that of the BSE run of \cref{app:E.15} with the exchange switched off, so it costs little on top of that run, and the limit without interaction, which must return a quarter of the independent-particle density response, controls it.

\subsection{Anisotropy of the electronic response}

A cubic crystal's response depends on the direction of $\vect{q}$, and that dependence is what turns the Earth's rotation into a daily modulation (\cref{sec:11.7}). It is carried by the cubic-harmonic components $S_\ell(|\vect{q}|,\omega) = \int\dd\hat{\vect{q}}\,Y^{\mathrm{c}}_\ell(\hat{\vect{q}})\,S(\vect{q},\omega)$ on the cubic harmonics $Y^{\mathrm{c}}_\ell$ of \cref{app:C.3}, for $\ell = 4$, 6 and 8, reported as $S_\ell/S_0$. Odd $\ell$ vanish by time reversal and $\ell = 2$ by cubic symmetry. The loss function with local fields is computed at full $\vect{q}$ (GPAW or BerkeleyGW) on a Lebedev grid of directions and projected. Near the gap the response has structure on a scale of \SI{0.1}{rad}: at $|\vect{q}| = 1.2$ to \SI{1.48}{keV} the zero-phonon edge of silicon has its minima 0.13 to \SI{0.22}{rad} from the valley axis in the $\{001\}$ planes, 14 to \SI{107}{meV} below the on-axis value, and shallower minima in the $\{011\}$ planes (\cref{sec:13.9}). A 110-point Lebedev rule integrates polynomials exactly only to degree 17, which is not enough once the response has components with $\ell \ge 10$. The grid is therefore refined from 110 to 230 to 434 points, exact to degrees 17, 25 and 35, and the test is that $S_4$, $S_6$ and $S_8$ change by less than \SI{1}{\percent} between the last two. Near the edge the adaptive cone grid of \cref{app:E.14} is added. Time reversal ($\ell = 1, 3$) and cubic symmetry ($\ell = 2$) supply checks that cost nothing: those components stay below $10^{-3}$ of the isotropic part. Projecting a line that is sharp in $\omega$ before binning gives meaningless coefficients, so $\omega$ is binned first. The run shares its band calculation with the library and adds little cost.

\subsection{Local fields in the phonon-assisted window}

The rate estimate of silicon in \cref{sec:13.9} is normalized to the measured optical edge and carries a screening correction relative to that calibration: 0.3 to 2.6 for the external factor and the local-field wings, an estimate, and the pair channel of \cref{sec:13.4} beyond the calibration, which no estimate sizes. A model, a supercell calculation and a bound from zero-phonon vertices size the correction.

The model tests the method itself. A one-dimensional crystal with two cosine bands, each of width 2 and separated by a gap of 1 in model units, a contact Coulomb kernel $v_0$ for wavevectors below half the zone, a classical phonon field $u_0\cos\omega_Qt$ at wavevector $Q$, and a probe at $(q,\omega)$ has a second-order induced density at the pair momentum $q+Q$ and frequency $\omega_f = \omega-\omega_Q$. It is computed twice: in the frequency domain, from the screened amplitude of \cref{sec:13.4}, and in real time, by propagating the density matrix in the time-dependent Hartree approximation with both perturbations and Fourier transforming. The limits of three couplings are known: a rigid coupling, $g = -u\,\partial_xV$, gives a screening ratio $\mathfrak{r}_{\mathrm{scr}} = |\tilde M|^2/|M|^2$ of order $\omega_Q/\omega$; a coupling proportional to the density vertex gives $\tilde M = M/\varepsilon(q+Q,\omega_f)$ exactly; a generic intracell coupling gives $\mathfrak{r}_{\mathrm{scr}}$ of order $1 - 1/\varepsilon(q+Q)$. The model uses $\omega_Q = 0.05$, $Q$ at a quarter of the zone, $u_0$ small enough that the response is linear in it (checked by halving), 2000 to 8000 $k$ points, $v_0$ scanned so that $\varepsilon(q+Q)$ runs from 1.5 to 4, a fourth-order Runge--Kutta step of $0.01/\Eg$, a total time of at least $200/\Eg$ and an adiabatic switch-on. The two routes agree to $10^{-3}$, stable under doubling the $k$ points and halving the step. On-shell residues cannot be compared on a finite ring, so the comparison uses many $k$ points and a damping extrapolated to zero. The model runs in minutes.

The silicon calculation uses a supercell. A supercell displaced along the special (zero-gradient) configuration of \cref{app:E.8} samples the phonon average deterministically, and the loss function of the displaced supercell, unfolded to the primitive momentum, contains the phonon-assisted response at the momenta the supercell resolves, and three runs of it separate the corrections. The supercell with RPA local fields in its own reciprocal-lattice basis gives the full result; the same supercell with independent particles and a scalar screening $1/\varepsilon(\vect{q},\omega)$ gives the reference, and the difference is $\delta_{\mathrm{LF}} + \delta_{\mathrm{pair}}$; a third run with local fields restricted to the primitive reciprocal lattice gives $\delta_{\mathrm{LF}}$ alone. On $\mathfrak{T}$ this is $n_B = 0$, $n_g = 2$, $\mathcal{X}$ = RPA with local fields; the special displacements are the adiabatic quadrature, and the correction of order $\omega_Q$ of \cref{sec:13.2} is applied afterwards. Cubic supercells of $n\times n\times n$ conventional cells resolve only certain momenta: $n = 3$ (216 atoms) reaches $|\vect{q}| = 0$, 0.76, 1.08, 1.32 and \SI{1.52}{keV}, and $n = 4$ (512 atoms) reaches 0, 0.57, 0.81, 0.99 and \SI{1.14}{keV}. The two meshes share no nonzero momentum below \SI{2.28}{keV}, so a convergence test between them is impossible at fixed $|\vect{q}|$. The tests are instead the change from $n = 2$ (64 atoms) to $n = 4$ at their common \SI{1.14}{keV}, below \SI{5}{\percent}, and the agreement of $n = 3$ and $n = 4$ after interpolation in $|\vect{q}|$ between 0.76 and \SI{1.52}{keV}, also within \SI{5}{\percent}. At $n = 2$ the special configuration samples only a $2^3$ phonon mesh, so the first test also measures phonon sampling; both antithetic configurations are reported. The ground state is PBE with a scissor; the screening uses Yambo or a Liouville--Lanczos solver, which scales better, with a local-field cutoff of \SI{3}{Ry} and bands to \SI{30}{eV} above the Fermi level, and $\mathrm{\Gamma}$-only sampling checked against $2^3$. The output is the screening ratio $\mathfrak{r}_{\mathrm{scr}}(\vect{q}) = 1 + \delta_{\mathrm{LF}} + \delta_{\mathrm{pair}}$, the full result over the reference, at the resolved momenta, and the test $\mathfrak{r}_{\mathrm{scr}}(\vect{q}) \approx \mathfrak{r}_{\mathrm{scr}}(0)$. The three cell sizes cost about $10^2$, $2\times10^3$ and $2\times10^4$ core-hours, hours to a day on accelerators. A second pass with the adiabatic local-density kernel isolates its extra source term, estimated at a few percent of $L^{\mathrm{PA}}$ (\cref{sec:13.4}).

The bound needs no amplitude. The external local-field correction alone can be bounded from the zero-phonon vertices at the same $(\vect{q},\omega)$,
\begin{equation}
  1 + \delta_{\mathrm{LF}}^{(0)}(\vect{q},\omega) = \frac{\sum_{cv}\big|\sum_{\vect{G}}(\tilde\varepsilon^{-1})_{\vect{G}0}(\vect{q},\omega)\,\rho_{cv}(\vect{q}+\vect{G})\big|^2}{\sum_{cv}|\rho_{cv}(\vect{q})/\varepsilon_{\mathrm{mac}}|^2} , \qquad \varepsilon_{\mathrm{mac}} = 1/(\tilde\varepsilon^{-1})_{00} ,
  \label{eq:E.16-dLF}
\end{equation}
a ratio that is one without local fields. It needs the column $(\tilde\varepsilon^{-1})_{\vect{G}0}(\vect{q},\omega)$ in the primitive cell for $\vect{q}$ on the library grid and $\omega$ from 1.2 to \SI{1.6}{eV}, with 40 bands, a $12^3$ grid and a cutoff of \SI{4}{Ry} (about 37 reciprocal-lattice vectors), converged to \SI{1}{\percent} between 3 and \SI{5}{Ry}. It costs about 10 core-hours.

\subsection{Loss functions of the candidate targets}

The ranking of targets in \cref{sec:21.7} needs the loss function $L(\vect{k},\omega) = \Imag[-(\varepsilon^{-1})_{00}(\vect{k},\omega)]$ of each candidate for $|\vect{k}|$ up to about \SI{5}{keV} and $\omega$ from the gap to about five valence plasma energies, 75 to \SI{150}{eV}. Its inverse fourth moment in $\omega$ sets the Migdal yield, and its weight near the gap sets the electron-recoil rate at low mass. The level is RPA with the full matrix $\varepsilon_{\vect{G}\vect{G}'}$, since $L$ is not $\Imag\varepsilon_{00}/|\varepsilon_{00}|^2$ once local fields matter, on a ground state with a scissor or quasiparticle gap matched to the measured one, with the adiabatic local-density kernel as the cross-check. A scissor serves only where PBE gets the band order right, checked against measured gaps at the high-symmetry points; elsewhere the bands must pass the check of order, gaps and masses of \cref{app:E.14}, as for germanium. Semicore shells within five plasma energies of the gap are included: Ga 3d, Ge 3d, W 5p and 4f, Pb 5d, Te 4d, Cs 5p, I 4d, Mo 4p. A cubic crystal needs three directions, a uniaxial one the directions parallel and perpendicular to the axis and one oblique. Every directional average is over $\hat{\vect{k}}$, never one direction.

Since the inverse fourth moment weights the gap as $\Eg^{-5}$, every grid test is run on that moment, not on $L$ itself. The $f$-sum rule $\int\omega L\,\dd\omega = (\pi/2)\omega_p^2$ holds to \SI{2}{\percent} with the electrons counted (valence plus the semicore included), which also fixes the split between valence and core. The number of bands, the $\vect{k}$ grid and the local-field cutoff are each increased until $\int L/\omega^4\,\dd\omega$ changes by less than \SI{1}{\percent}. The broadening is at most $\Eg/20$ and is extrapolated to zero, since the inverse fourth moment is weighted toward the gap. Missing semicore weight biases that moment low at a fixed $f$-sum, because the weight then sits at high $\omega$. The scissor moves the onset, which the moment weights as $\Eg^{-5}$, and the computed moment is compared with the bound $(\pi/2)\omega_p^2/\Eg^5$ that the $f$-sum rule gives. Excitonic weight below the gap, absent in the RPA, adds to the moment; the BSE of \cref{app:E.15} bounds it. Beyond the first zone, the element at $\vect{G} = 0$ must be taken at the unreduced $\vect{k}$ (\cref{app:E.6}). We estimate the cost from \cref{we:E-elf-si} at tens to a hundred core-hours for a two-atom cell and several hundred for a ten-atom oxide.

\subsection{Relativistic shift of the \texorpdfstring{$f$}{f}-sum rule}

The $f$-sum rule holds exactly only for a nonrelativistic Hamiltonian. At order $\alpha^2$, the mass--velocity term reduces it, by $1.8\times10^{-3}$ in silicon (\cref{sec:4.3}), and so does the orbit--orbit (Breit) interaction between electron pairs, which \cref{sec:10.6} writes as an integral of the all-electron static structure factor over the Coulomb kernel with a transverse projector. The all-electron atoms of \cref{sec:10.6} put the total relativistic shift at small $q$ at $-0.24\,\%$ for silicon and $-0.76\,\%$ for germanium, of which the mass--velocity part is $-0.18\,\%$ and $-0.58\,\%$; a Thomas--Fermi estimate gives $-0.23$ and $-0.73\,\%$. What is missing is the orbit--orbit shift as a function of $q$ from 1 to \SI{100}{keV}, and the total with the mass--velocity term. The cores dominate it, so the calculation is atomic: Hartree--Fock with the exact two-component (X2C) scalar-relativistic Hamiltonian in pyscf, in an uncontracted large basis (the atomic natural orbital (ANO) set for relativistic core correlation (RCC), ANO-RCC, or cc-pwCV5Z) with tight $s$ and $p$ exponents up to $10^7$, gives the two-electron energy and, from the one-particle density matrix in momentum space, the static structure factor. A crystal correction for the valence comes from a PAW all-electron density with the electron-gas exchange--correlation hole. The two-electron energy is converged to $10^{-4}$ when the basis goes from quadruple to quintuple zeta, and the $k$ integral to $10^{-3}$ with a radial grid out to $50/a_0$ for the germanium 1s shell; contracted valence basis sets would underestimate the core cusp by several percent. The orbit--orbit term is about 0.31 of the mass--velocity one in both atoms, so including only one misstates the total by about \SI{24}{\percent}. The formula subtracts the uniform Hartree term itself, so it is never subtracted by hand, and for the valence, where $q$ times the pair distance exceeds one, the full angular dependence is kept rather than its average. The atoms cost minutes, the crystal correction about 10 core-hours, and the result replaces the estimate in \cref{sec:10.6}.

\subsection{Static screening with lattice relaxation}

A probe that couples only to electrons still moves the nuclei in a static measurement, because the electrons drag them. The static response of the electron channel with the lattice allowed to relax is the clamped-nuclei step plus a lattice term (\cref{sec:10.2}),
\begin{align}
  1 - \varepsilon_e^{-1}(\vect{q},0) &= [1 - \varepsilon^{-1}_{\mathrm{cl}}(\vect{q})] + \mathcal{G}_{\mathrm{rel}}(\vect{q}) , \nonumber\\
  \mathcal{G}_{\mathrm{rel}}(\vect{q}) &= \frac{e^2}{q^2\Omega}\sum_\nu\frac{\big|\sum_\kappa\vect{q}\cdot Z^{(e)}_\kappa(\vect{q})\cdot\vect{e}_{\kappa\nu}(\vect{q}_r)\,\ee^{\ii\vect{q}\cdot\vect{\tau}_\kappa}/\sqrt{M_\kappa}\big|^2}{\omega_\nu(\vect{q}_r)^2} ,
  \label{eq:E.16-step}
\end{align}
at adiabatic, harmonic, one-phonon order. Here $Z^{(e)}_\kappa(\vect{q})$ is the electron charge that follows atom $\kappa$ when it moves, which at $q\to0$ in an insulator is $Z_\kappa - Z^*_\kappa$ with all electrons counted, 14 for silicon. As $q\to0$ the acoustic branch dominates and $\mathcal{G}_{\mathrm{rel}} \to \varOmega_{P,e}^2/(c_{\mathrm{LA}}q)^2$, with $\varOmega_{P,e}^2 = e^2\nd_e^2/\rho_T$ and $c_{\mathrm{LA}}$ the longitudinal sound speed: the step diverges as $1/q^2$ in every solid. For silicon $\varOmega_{P,e} = \SI{0.51}{eV}$, and the limiting form reaches one at \SI{18}{keV}, above the zone boundary; it gives $3\times10^2$ at \SI{1}{keV}, where neither the linear dispersion nor $Z^{(e)} \approx 14$ holds, so the true value there is only estimated at $\gtrsim10^2$.

The calculation needs $Z^{(e)}_\kappa(\vect{q})$ at finite $q$. Frozen-phonon supercells commensurate with $\vect{q}$, for silicon $1\times1\times N$ conventional cubic cells along [100] with $N = 2$ to 8, displace one sublattice by $u\,\ee^{\ii\vect{q}\cdot\vect{R}}$; the Fourier component of the induced electron density at $\vect{q}$, at two amplitudes to check linearity, gives the charge. The DFPT induced density is equivalent. A pseudopotential calculation gives only the valence part; the frozen cores are added back as rigid charges, valid for $q$ well below the inverse core radius, about \SI{30}{keV} for the silicon 1s shell. The phonons come from \cref{app:E.17}, the clamped step from the loss functions above. $Z^{(e)}$ has to be stable to \SI{1}{\percent} in cutoff and $\vect{k}$ grid and $\mathcal{G}_{\mathrm{rel}}$ to \SI{2}{\percent} under refinement in $q$, and the small-$q$ limit has to match $c_{\mathrm{LA}}$ from the computed elastic constants. A violated acoustic sum rule sends $\omega_\nu \to 0$ with a finite numerator and fakes a divergence, so the sum rule is enforced on the force constants. The momentum is kept unreduced, $\vect{q} = \vect{q}_r + \vect{G}$, so that the step follows the density to several zone momenta, with its Umklapp terms; above a Debye--Waller exponent $2W \approx 0.1$ multiphonon terms are added; and in GaAs the direction dependence of the sound speed and of the nonanalytic optical term must be kept. Silicon and GaAs together cost about $10^2$ core-hours.

\subsection{Kinetic-energy response}

A scalar coupled to the electron's mass is absorbed through the kinetic-energy density $\mathcal{T}_e$, not the density (\cref{sec:10.7}). For every crystal whose point group has no polar vector, the absorption rate is proportional to $-\Imag R_{\mathcal{T}\mathcal{T}}(0,m_\phi)$ (\cref{eq:10.7-scalar}), whose independent-particle form is a joint density of states weighted by second-order momentum matrix elements; with $\mathcal{T}_e = \sum_jp_j^2/2m_e$ as in \cref{sec:8.1}, the squared element carries $1/4m_e^2$,
\begin{equation}
  \Imag R^0_{\mathcal{T}\mathcal{T}}(0,\omega) = -\frac{2\pi}{\Vol}\sum_{cv\vect{k}}\frac{|\langle c\vect{k}|p^2|v\vect{k}\rangle|^2}{4m_e^2}\,\delta(\omega - \varepsilon_{c\vect{k}} + \varepsilon_{v\vect{k}}) ,
  \label{eq:E.16-TT}
\end{equation}
with spin summed. The matrix element is computed directly from the plane-wave coefficients,
\begin{equation*}
  \langle c\vect{k}|p^2|v\vect{k}\rangle = \sum_{\vect{G}}|\vect{k}+\vect{G}|^2\,c^*_{c\vect{k}}(\vect{G})\,c_{v\vect{k}}(\vect{G}) ,
\end{equation*}
plus the PAW one-center correction. It equals $\sum_m\vect{p}_{cm}\cdot\vect{p}_{mv}$ over all bands including the core, a sum that converges slowly. The kinetic matrix element converges like the total kinetic energy, so silicon needs at least \SI{80}{Ry} with SG15; the grid needs at least $24^3$ points, the delta function a broadening of at most \SI{20}{meV}, and the gap a scissor. The $f$-sum rule is not a check here. Its analogue for $\mathcal{T}_e$ fixes the energy-weighted integral by $\langle|\nabla V|^2\rangle/2m_e$ per electron (\cref{sec:10.7}), which the core region dominates, and the check uses that sum rule restricted to the valence. For the same reason, norm-conserving pseudo-wavefunctions, which lack the core oscillations, are compared with a PAW or all-electron calculation. $\mathcal{T}_e$ is the canonical kinetic operator; velocity matrix elements $\ii[H,\vect{r}]$ are not a substitute. Since $\mathcal{T}_e$ is even under inversion, the direct edge of silicon at $\mathrm{\Gamma}$ is suppressed, and absorption inside the gap needs the phonon-assisted row of \cref{app:E.14}. Polar crystals add a mixing term through the current response. Local fields, the next many-body approximation on $\mathfrak{T}$, are added at RPA level. Silicon takes minutes; a ten-atom oxide about one accelerator-hour.

Operator $\mathcal{O}_2$ needs the same row at finite momentum. Its four functions per version, $\mathcal{T}_e(\vect{q})$ with itself and with the density and the two transverse currents, enter the Dyson solve of the library as external rows and columns, like the current, so that they are screened by the same density block (\cref{sec:21.6}). The vertex $\langle c,\vect{k}+\vect{q}|\mathcal{T}_e(\vect{q})|v\vect{k}\rangle$ is built from the plane-wave coefficients with the matrix element $(\vect{k}+\vect{k}')^2/8m_e$ between plane waves of \cref{eq:8.1-rows}, plus the PAW one-center correction, on the grid of the library, and as $\vect{q}\to0$ it must return \cref{eq:E.16-TT}.

\subsection{Cosmic-ray-boosted dark matter}

For dark matter boosted to relativistic speeds, the density row of the response is contracted with the full current. The published treatment of cosmic-ray-boosted dark matter scattering in silicon and germanium~\cite{liang2025plasmon} drops the factor $(1 - \omega^2/\mathcal{Q}^2)^2$ of the exact density row, with $\mathcal{Q}$ the magnitude of the momentum transfer, and the transverse row (\cref{sec:10.9}). Since $\omega \le v_\chi\mathcal{Q}$, the rate changes by at most \SI{2}{\percent} for $v_\chi \lesssim 0.1$, but for a cosmic-ray-boosted flux with $v_\chi \to 1$ the bound is empty, and a transverse row, $(v_\chi^2 - \omega^2/\mathcal{Q}^2)\,\Imag[-\mathcal{Q}^2/(\omega^2\varepsilon^{\mathrm{T}} - \mathcal{Q}^2)]$ per unit of the density normalization, adds to it. The recast re-implements the published rate with the exact density row and the transverse term: the source's random-phase longitudinal function, a transverse function equal to it as $\mathcal{Q}\to0$ and of Lindhard form otherwise, the source's flux tables, and a three-dimensional quadrature in $\mathcal{Q}$, $\omega$ and the flux angle. The first step is to reproduce the published limit with the published rate to \SI{5}{\percent}. The transverse term has a Cherenkov-type pole where $\omega^2\Real\varepsilon^{\mathrm{T}} = \mathcal{Q}^2$, which needs a scan in the imaginary part of the frequency, and the mediator form factor is evaluated at $\mathcal{Q}^2 - \omega^2$, not at $\mathcal{Q}^2$. For a heavy mediator the exact density row changes the rate by $10^{-4}$ and the transverse row adds at most $10^{-4}$, so that limit is unchanged; since the boosted flux is itself proportional to $\sigmae$, a change $R$ of the detection rate moves a limit by $R^{-1/2}$ (\cref{sec:21.5}). For a massless mediator the shift of $\sigmae$ is below one percent for $\mchi \ge \SI{10}{keV}$, and below about \SI{3}{keV} it is undetermined: the published density row has a pole at $\mathcal{Q} = \omega$, whose weight depends on the treatment of the lowest kinetic energies of the flux, which the source does not state. Depending on that treatment the ratio of the density rows ranges from 0.05 to 0.93 at \SI{1}{keV}, so the recast must state the treatment it adopts, and its light-mediator result stands only with that statement. The recast costs a few core-hours.
